% !TEX root = ../nc-quadric-avatar.tex
% 2026-11-23 / 相対square完成と実際のtails認識。
\section{相対square完成のAS解消とtails認識}
\label{sec:qa-relative-square-completion}

\subsection{平坦compact対象の相対Ext次元}
\label{subsec:qa-flat-compact-ext-dimension}

以下では基底を局所有限型 $k$-scheme とし，必要に応じて共通の fppf
局所 affine chart $S=\operatorname{Spec}R$ を取る．特異な基底上で
coherent heart 全体の大域次元が二であるとは仮定しない．

\begin{lemma}[平坦compactな始域に対する相対Serre評価]
\label{lem:qa-flat-compact-ext-two}
$\mathcal C$ を相対 $H_2$ の heart とし，その局所 compact 圏が
$R$ 上有限 projective で相対 smooth な代数 $\Gamma$ によって
Morita 表示されているとする．$E\in\mathcal C$ が $R$ 上平坦かつ
compact なら，任意の $K\in\operatorname{Ind}\mathcal C$ に対して
\[
 \operatorname{Ext}^q_{\operatorname{Ind}\mathcal C}(E,K)=0
 \qquad(q>2)
 \tag{FE2}\label{eq:qa-flat-compact-ext-two}
\]
が成立する．この結論は同じ仮定を満たす任意基底変更後にも成立する．
\end{lemma}
\begin{proof}
$I$ を $R$-加群圏の injective cogenerator とする．
相対 Serre 双対を compact な第二変数だけに限定せず，次の随伴として使う：
\[
 \mathbf R\operatorname{Hom}_R
   (\mathbf R\operatorname{Hom}_{\mathcal C}(E,K),I)
 \simeq
 \mathbf R\operatorname{Hom}_{\mathcal C}
   (K,\mathsf S E\otimes_R^{\mathbf L}I).
 \tag{SA}\label{eq:qa-serre-injective-adjunction}
\]
この式は任意の $K\in D(\operatorname{Ind}\mathcal C)$ に対して成立する．
実際，Morita 表示で $E$ に対応する perfect 右 $\Gamma$-加群を $Q$，
\[
 Q^\vee=\mathbf R\operatorname{Hom}_\Gamma(Q,\Gamma)
\]
をその左双対とする．$\Gamma$ が $R$ 上 perfect なので $Q^\vee$ も
$R$ 上 perfect である．tensor--Hom 随伴と有限 perfect 双対によって
\[
 \begin{split}
 \mathbf R\operatorname{Hom}_R
  (K\otimes_\Gamma^{\mathbf L}Q^\vee,I)
 &\simeq \mathbf R\operatorname{Hom}_\Gamma
     (K,\mathbf R\operatorname{Hom}_R(Q^\vee,I))\\
 &\simeq \mathbf R\operatorname{Hom}_\Gamma
     (K,\mathbf R\operatorname{Hom}_R(Q^\vee,R)
                    \otimes_R^{\mathbf L}I)
 \end{split}
\]
を得る．最後の右加群は $\mathsf S Q\otimes_R^{\mathbf L}I$ であり，
これが式~\eqref{eq:qa-serre-injective-adjunction} である．
ここで $K$ の compactness や基底平坦性は使っていない．

相対 $H_2$ では $\mathsf S E=t^{-1}E[2]$ であり，canonical $t$ は
heart 上完全で $R$-線形である．従って $t^{-1}E$ も基底平坦で，
\[
 \mathsf S E\otimes_R^{\mathbf L}I
       =(t^{-1}E\otimes_R I)[2]
\]
となる．括弧内は Ind heart の対象である．$K$ も heart にあるので，
式~\eqref{eq:qa-serre-injective-adjunction} の右辺は cohomological
次数 $-2$ 未満で零である．一方 $I$ の injectivity により，左辺の
次数 $-q$ cohomology は
\[
 \operatorname{Hom}_R
       (\operatorname{Ext}^q_{\mathcal C}(E,K),I)
\]
である．$I$ が cogenerator であることから
式~\eqref{eq:qa-flat-compact-ext-two} が従う．
\end{proof}

\begin{proposition}[有限個の平坦対象を加えた相対AZ生成]
\label{prop:qa-relative-flat-foundation-az}
$\mathcal C,\mathcal O$ を相対 $H_2$ 族とし，
$E_0=\mathcal O,E_1,\ldots,E_r$ を基底平坦 compact な heart 対象とする．
各 affine 基底変更の共通 fppf 局所 chart 上で前補題の有限 Morita
表示があり，そこで各 $i$ に対し
\[
 \operatorname{Ext}^q(E_i,t^m\mathcal O)=0
 \qquad(q>0,\ m\gg0)
 \tag{FO}\label{eq:qa-flat-foundation-forward-o}
\]
を仮定する．このとき
\[
 T=\bigoplus_{i=0}^r E_i
\]
に対し $(T,t)$ は各 affine 基底変更上で AZ ample である．
\end{proposition}
\begin{proof}
$(\mathcal O,t)$ は論文Bの相対反 Serre 再構成により AZ ample であり，
$\mathcal O$ は $T$ の直和因子なので生成条件 AZ1 は直ちに従う．

まず上記の局所 Morita chart 上で議論する．AZ2 のため，任意の
coherent 対象 $K$ に対して
\[
 \operatorname{Ext}^1(t^{-n}E_i,K)=0\quad(n\gg0)
 \tag{FV}\label{eq:qa-flat-foundation-ext-one-vanish}
\]
を示す．AZ1 を二回使い，短完全列
\[
 0\longrightarrow L\longrightarrow P\longrightarrow K\longrightarrow0,
 \qquad
 0\longrightarrow L'\longrightarrow P'\longrightarrow L\longrightarrow0
\]
を取る．ここで $P,P'$ は有限個の $t^{-a}\mathcal O$ の直和である．
基底の正則性も $K,L,L'$ の基底平坦性も仮定しない．
条件~\eqref{eq:qa-flat-foundation-forward-o} により，十分大きい $n$ について
\[
 \operatorname{Ext}^1(t^{-n}E_i,P)=0,
 \qquad \operatorname{Ext}^2(t^{-n}E_i,P')=0
\]
となる．さらに $t^{-n}E_i$ は基底平坦 compact なので，前補題から
$\operatorname{Ext}^3(t^{-n}E_i,L')=0$ である．
二つの長完全列により
\[
 \operatorname{Ext}^1(t^{-n}E_i,K)
  \hookrightarrow\operatorname{Ext}^2(t^{-n}E_i,L)=0
\]
を得る．これが式~\eqref{eq:qa-flat-foundation-ext-one-vanish} である．
任意の全射の kernel にこの消滅を適用すれば AZ2 を得る．
各 $i$ は有限個なので，一つの十分大きい $n$ を選べる．

一般の affine noetherian 基底 $\operatorname{Spec}R$ 上の AZ2 を確認する．
その準compactnessにより，共通 Morita chart となる有限個の affine fppf
被覆を選べる．その非交和を $\operatorname{Spec}R'$ とすれば，
$R\to R'$ は faithfully flat である．与えられた全射を引き戻すと，
有限個の chart 上の消滅から一つの十分大きい $n$ の範囲で
\[
 \operatorname{Hom}(t^{-n}T,M)\otimes_R R'
 \longrightarrow
 \operatorname{Hom}(t^{-n}T,N)\otimes_R R'
\]
が全射となる．ここでは相対 $H_2$ の coherent Hom の flat base change
と $t$ の基底変更整合性を用いている．余核を faithfully flat に
引き戻すと零なので，もとの Hom 写像も全射である．
これにより AZ2 は基底全体で成立する．この最後の議論には大域の
$D(\operatorname{Ind}\mathcal C)$ に対する Morita 表示を使っていない．
\end{proof}

\begin{remark}[AZ生成における構造対象の位置と二段presentation]
\label{rem:qa-relative-az-noninitial-structure}
命題~\ref{prop:qa-relative-flat-foundation-az}では，
指定構造対象$\mathcal O$が$T$のいずれかの
直和因子であればよく，先頭の因子である必要はない．
実際，AZ1は$\mathcal O$の直和因子性のみを使い，
AZ2の証明は各$E_i$から十分前向きの
$t^m\mathcal O$への高次Ext消滅のみを使う．
二回の$\mathcal O$-軌道対象からの表示
$0\to L\to P\to K\to0$，
$0\to L'\to P'\to L\to0$は
$\operatorname{Ext}^1(E,P)=
 \operatorname{Ext}^2(E,P')=
 \operatorname{Ext}^3(E,L')=0$
から$\operatorname{Ext}^1(E,K)=0$を与える．
最後の消滅には平坦compact始域の
補題~\ref{lem:qa-flat-compact-ext-two}を使い，
基底の絶対大域次元二は仮定していない．
従ってこれは反対normalの第一頂点が
強い$H_2$ pointingでなくても，
他の頂点に適合構造対象を含む場合に適用できる．
ただし全ての候補対象が\emph{同じ相対heart}の
基底平坦compact対象であることは別途確認しなければならない．
\end{remark}

\subsection{相対square完成の標準tails心}
\label{subsec:qa-relative-square-tails-heart}

\begin{theorem}[square完成の相対標識付きtails認識]
\label{thm:qa-relative-square-completion-tails}
相対 curve-type square avatar の実際の有限代数を $\Lambda$，
その四つの標準対象を $E_0,\ldots,E_3$ とする．
定理~\ref{thm:qa-relative-square-e0-e3-h2} の相対 $H_2$ heart
$\mathcal C_0$ を用い，$T=\bigoplus_{i=0}^3E_i$ と置く．
その通常の相対 Serre 完成を
\[
 C(\Lambda)=\bigoplus_{n\geq0}
 H^0(\Theta_\Lambda^{\otimes_\Lambda^{\mathbf L}n}),
 \qquad
 \Theta_\Lambda=\mathbf R\operatorname{Hom}_{\Lambda^e}
                      (\Lambda,\Lambda^e)[2]
\]
とする．全ての本稿の affine 基底変更上で，
$C(\Lambda)$ は右 noetherian であり，実際の完全同値
\[
 \qgr C(\Lambda)\xrightarrow{\sim}\mathcal C_0
 \tag{SQH}\label{eq:qa-relative-square-completion-heart}
\]
が存在する．標準冪等元を $e_i$ とすると，この同値は
$\pi(e_iC(\Lambda))\mapsto E_i$ と次数移動 $(1)\mapsto t$ を保つ．
四頂点の周期的 unfolding による square $\mathbb Z$-代数を
$A(\Lambda)$ と書けば，同じ主張は $\qgr A(\Lambda)$ にも成立する．

同値~\eqref{eq:qa-relative-square-completion-heart} は任意の本稿の
基底変更に対する module-object 引戻しと coherent に整合する．
次数別代数 $C(\Lambda)$，$A(\Lambda)$ の形成自体は，前節の
Serre tensor 比較により任意の scheme 基底変更と可換する．
\end{theorem}
\begin{proof}
前節で全 $E_i$ は $\mathcal C_0$ の基底平坦 compact な heart 対象であり，
$T$ の有限 End は実際の $\Lambda$ と同定されている．
全 forward Hom 複体は基底上 perfect で，全幾何 fiber の geometric helix
上の前向き Ext 消滅から次数零に集中する．特に $m\geq1$ について
\[
 \mathbf R\operatorname{Hom}(E_i,t^mE_0)
\]
は次数零の有限局所自由層である．これは任意基底変更後にも成り立つ．
従って前命題の条件~\eqref{eq:qa-flat-foundation-forward-o} が成立し，
$(T,t)$ は相対 AZ ample である．

通常の軌道環
\[
 R(T,t)=\bigoplus_{n\geq0}\operatorname{Hom}(T,t^nT)
\]
に対し，$\operatorname{End}(T)=\Lambda$ は右 noetherian であり，
全 coherent Hom は有限 $R$-加群なので有限 $\Lambda$-加群でもある．
一般の End 環に対する Artin--Zhang 再構成を適用すると，
$R(T,t)$ の右 noetherian 性と
\[
 \mathcal C_0\simeq\qgr R(T,t)
\]
を得る．ここで用いる形は，論文B
\cite[\texttt{prop:relative-h2-antiserre-reconstruction}]{QuadricModuli2026}
が引用する \cite[Theorem 4.5(1)]{AZ1994}，
\cite[Theorem 7.3]{OACU2026} と同じである．

共通の fppf 局所 chart の実際の tilting Morita 同値の下で $t$ は
$-\otimes_\Lambda^{\mathbf L}\Theta_\Lambda$ に対応する．
従って
\[
 \mathbf R\operatorname{Hom}(T,t^nT)
  \simeq \Theta_\Lambda^{\otimes_\Lambda^{\mathbf L}n}
\]
は積を保つ同型である．全次数の次数零集中を取り，coherent Hom の
flat base change に沿って降下すると，基底全体で
$R(T,t)=C(\Lambda)$ を得る．四つの冪等元による成分分解は標準
$E_i$ を保つ．有限個の頂点ごとに graded 成分を並べ直す通常の
folding/unfolding は graded 加群圏の同値を与え，有限生成性と
torsion を保つので $\qgr A(\Lambda)$ についても同じ結論となる．

最後に，AZ 同値の逆向きは graded 自由加群のシフトを
$t^{-n}T$ に送り，有限表示の余核へ延長して作られる．
通常の module-object 引戻しは右完全なので，この有限表示の構成と
可換する．標準射影 $e_iC(\Lambda)$ についても同様である．
自由対象上の比較は合成に対して coherent であり，有限表示への
延長の一意性から全対象上の比較も coherent となる．
非平坦基底変更に対して全対象の低次数 $\Gamma_*$ が可換すると
仮定する必要はない．
非 affine な基底ではこの比較に沿って Abel 圏の descent を適用し，
実際の完全同値を貼り合わせる．
大域の $D(\operatorname{Ind}\mathcal C_0)$ と局所 compact 圏の descent
を同一視する追加命題は用いていない．
\end{proof}

\begin{remark}[結論の範囲]
これは実際の曲線 avatar 族から構成した有限 square 代数とその canonical
Serre 完成に対する認識である．任意の AS/QP 族の Noetherian 性，任意の
presentation の相対 $H_2$ 性，あるいは全 presentation stack と
$\mathfrak H_2$ の同一視を結論していない．また
式~\eqref{eq:qa-flat-compact-ext-two} は始域を基底平坦 compact とした
相対的な消滅であり，特異な基底上の heart 全体の大域次元二という
主張ではない．
\end{remark}
